\documentclass[12pt,a4paper]{article}
\usepackage[UTF8]{ctex}
\usepackage{geometry}
\geometry{left=2.5cm,right=2.5cm,top=2.5cm,bottom=2.5cm}
\usepackage{enumitem}
\usepackage{amsmath}
\usepackage{amssymb}
\usepackage{titlesec}
\titleformat{\section}{\Large\bfseries}{\thesection}{1em}{}
\titleformat{\subsection}{\large\bfseries}{\thesubsection}{1em}{}
\setlist{leftmargin=*,itemsep=0.1em,parsep=0.1em}

\title{\textbf{腾讯地图算法岗HR面试常见问题}}
\author{面试准备笔记}
\date{\today}

\begin{document}

\maketitle

\section{自我介绍}
面试官你好，我是黄思铭，目前是北京航空航天大学硕士在读，专业是应用数学，课题组方向是强化学习，同时我对推荐系统也深入研究。

在技术方面，我熟练掌握Python、Pytorch，同时对MATLAB和C语言也较为熟悉；在理论知识方面，我本研期间数分分析、高等代数、数理统计和深度学习等理论课程都能获得平均95甚至更高的分数，也多次获得了校级学业奖学金；在工程实现方面，我有较强的工程实现能力，例如简历里面提到的三个项目，分别是基于元学习的冷启动框架、基于CVaR正则的多智能体算法以及卫星对抗多智能体环境的设计。

因此，我希望把个人的技术放到真实业务场景中验证，尤其是参与地图上下点推荐或补贴算法上面。

\section{优点}
我觉得优点的话，我个人觉得有这么几个：

其一，我学习能力算是一个比较出众的优点，抛开做题不谈，我在面对一个新方向或者新任务时，通常能够比较快地建立整体框架，然后通过阅读文献、代码以及实践去补齐细节，实现入门。例如刚接触强化学习时候，都能够以较快的速度学习到一个能够应用的阶段。所以如果进入新的业务团队，我也有信心比较快地理解业务逻辑和技术栈。

其二，一旦定好目标，我的自驱力会很强。旦我明确了要解决的问题，我会主动拆解任务，给自己设定阶段目标，并持续推进，而不是等别人一步一步安排。比如我定好目标，准备推荐算法面试时候，我能够逐步拆解这一目标，从传统推荐各链路经典算法，再到代码复现，再到强化学习落地传统推荐，再到生成式推荐，一步一个脚印，持续不断推进。

其三，我很注重团队的关系，擅长沟通协作。在研究生做科研阶段，我会主动同步自己的进展和问题，与同门进行沟通，一起合作解决问题。我个人也很喜欢那种合作的氛围，而不是单兵作战的环境。

\section{缺点}
缺点的话，我觉得有一个缺点确实我目前做得不够好的地方，也是被导师多次批评的地方。就是我做一件事，总是喜欢把他的理论基础完全吃透，再去应用，也就是我不敢在没有十足把握情况下去进行一个事情。导师就跟我说，你这就是数学思维太深入了，AI这个行业，等你理论看完，别人都推进到不知道哪里去了，应该要边应用边去学习那些在应用里面涉及到的，你不懂的理论。

还有一个就是我可能需要一个mentor来指导，才能事半功倍。

\section{为什么选择腾讯}
因为俗话说得好，有鹅选鹅。不开玩笑，具体而言我认识有这么三点：
第一，技术平台方面：腾讯有非常丰富的真实用户场景，包括腾讯地图在内的各种产品，例如视频号、游戏、微信等，都能够直接服务大规模用户，而不是停留在离线指标上面。

第二，业务场景方面：以腾讯地图为例，对上下车点以及附近酒店进行推荐，还有补贴算法，都是比较契合我的技术兴趣。

第三，人才培养方面：腾讯有完善的实习生培养体系，对实习生的培养是舍得下功夫的，这一点我不仅在各大论坛上能看到，在身边已经在腾讯实习或者工作的人也会如此评价。

\section{职业规划}
我的职业规划大致可以分为三步：
第一阶段：夯实基础，深入业务。争取在短时间内，掌握组里的推荐系统全链路或者某个链路，能够独立负责具体业务场景的算法迭代。

第二阶段：强化学习落地。我会在夯实基础的前提下，尝试前沿领域中强化学习落地推荐系统的各种思路，依次来解决推荐中“探索-利用”和“长期价值”这两个平时很难完善的点。

第三阶段：生成式推荐。众所周知，大模型对整个算法领域冲击很大，生成式推荐很有可能将来会大面积替代传统推荐，必须跟上大模型的东风，不断关注前沿生成式推荐的发展，并尝试落地到当前业务中去。

\section{手里有几个offer}
我手里暂时有的offer只有个别中厂，还有个别大厂推进到较后面的步骤，但无法确定能拿到offer。但无论如何，对我而言，腾讯还是我的第一选择，因为腾讯本身的平台、以及我目前面试的业务以及对实习生的培养，我确实都是比较向往的，也很想接受腾讯的培训。

\section{坚持最长的一件事}
我坚持最长的事情，应该是读史。从小时候开始，我就很喜欢阅读历史，当然小时候都看的是历史小说，例如《三国演义》、《说岳全传》、《杨家将》等。长大后逐渐开始阅读正史，不过我一般不会直接读《二十四史》原文，因为确实比较枯燥，更多会看一些 B 站 UP 主的历史解说，这样更容易理解，也更容易坚持。

我觉得读史的意义不只是知道发生了什么，而是去理解古人在当时环境下为什么会做出某些选择。比如他面对的局势是什么，手里有哪些资源，受到哪些限制。然后我也会想，如果换成我处在类似情境下，我会怎么做，怎样选择才更合理。所以读史对我来说既是一种兴趣，也是一种思维训练，能让我更习惯从背景、条件和人物动机等多个角度去分析问题。

\section{遇到困难怎么解决？}
我觉得最重要的一点在于定位问题，困难究竟是哪一个环节出现的，所谓磨刀不误砍柴工，我认为先定位问题比盲目去寻求解决更好。一方面，定位问题之后，可以缩小解决方案的范围，加速困难的解决；另一方面，一旦定位问题，就能为以后碰到类似的问题提供经验教训，学会解决该类问题的通性通法。

当然定位问题之后，怎么解决这个问题也是很困难的一个步骤。通常又要拆分为调研解决方案、尝试不同解决方案、验证问题是否解决。


\section{一件难忘的事情（最成功，最有成就感）}
我印象比较深的一件事，是之前做军方项目时，算法终于跑出来的那一刻。其实就是简历里面第二个项目，基于CVaR正则的多智能体强化学习算法。

因为这是我读研后接触的第一个项目，个人
前期其实遇到了不少问题，比如不确定性量化怎么应用到多智能体上面？如何定位代码bug？模型效果不好是代码问题还是算法问题？那段时间我反复检查数据处理、模型结构、训练参数和评价指标，也和同学、老师甚至大模型高频讨论。

所以当最后实验真正跑通，并且结果符合预期的时候，我感觉特别有成就感。它不是简单地“代码能运行”，而是说明前面很多分析、尝试和调整是有效的。

这件事也让我意识到，做算法项目需要耐心和持续投入，不能因为一开始结果不好就放弃。这个过程也让我更加确定自己对算法方向是有兴趣的。


\section{兴趣爱好}
主要是读史、观影和看球赛。

我平时比较喜欢了解古代历史，尤其是通读各个朝代的发展脉络。不过我一般不会直接读《二十四史》原文，因为确实比较枯燥，更多会看一些 B 站 UP 主的历史解说，这样更容易理解，也更容易坚持。我觉得读史的意义不只是知道发生了什么，而是去理解古人在当时环境下为什么会做出某些选择。比如他面对的局势是什么，手里有哪些资源，受到哪些限制。然后我也会想，如果换成我处在类似情境下，我会怎么做，怎样选择才更合理。所以读史对我来说既是一种兴趣，也是一种思维训练，能让我更习惯从背景、条件和人物动机等多个角度去分析问题。

而看电影和看球赛相对而言就轻松了很多，主要是放松身心，让自己没有那么紧绷。


\section{在训练过程中，遇到最大的困难是什么？怎么解决的？}
\subsection{冷启动}
项目中最大的挑战是损失函数的设计。最初我按照经典 MAML 只用第二段 loss 来训练，但发现当新 item 完全没有交互数据时，推荐效果根本没好转。分析后发现，经典 MAML 假设新任务至少有少量样本可以进行梯度更新，但推荐系统的冷启动更极端——新 item 可能在最初几小时内都是零数据，根本没法做梯度更新，只能直接用生成的 embedding 预测。

为了解决这个问题，我改进了损失函数，加入了第一段 loss。第一段 loss 约束生成器输出的初始 embedding 本身就要好用（应对零数据场景），第二段 loss 约束更新后的 embedding 要更好（应对少量数据场景）。这样生成器需要同时优化初始质量和快速适应能力两个目标。
\subsection{强化学习}
这是军方的一个项目，当时一个核心要求是“不确定性量化”。而对于我而言，我之前根本没有接触过这个概念，其实不只是我，我们整个课题组都是这样。所以对于刚研一的我而言，这真是我入学以来遇到过最难的问题了。最终怎么解决这个问题呢？其实靠的是团结的力量，我们选择每周开一次不确定性量化讨论班。课题组所有人一起看一本不确定性量化相关书籍，每个人一周讲一个章节，一起讨论，遇到困难一起克服，最终把不确定性量化这种思想应用到了强化学习里面。
\subsection{卫星对抗}
我遇到的最大困难是奖励函数设计不好导致智能体学会了消极策略。具体表现是训练初期智能体不敢接近敌人，总是往战场边缘跑，甚至学会了在边界附近徘徊拖时间。这完全违背了我们希望它主动作战的初衷。

问题分析：我分析后发现主要有三个原因。第一是攻击敌人的奖励相比被敌人攻击的惩罚没有优势，智能体算了一笔账发现"进攻可能被打"和"进攻打中敌人"的收益差不多，那还不如不动。第二是没有范围惩罚，智能体可以无限远离战场中心，只要不出界就没事，所以它学会了躲在角落里。第三是缺少时间成本，智能体可以无限拖延，最后学会了拖到300步超时的策略。

在我不断改reward，做实验这样重复的摸索中，终于是设计出了三个额外的奖励，以使得最终成功：分别是探测奖励、时长惩罚和奖励分层
\end{document}